\section{Construction9: a charged Krom cut and two non-transitions}
\label{sec:construction9-main4}

The current achievable-region argument gains one exact partial-domain
point, not an all-input transition.  \textsc{STRONG-KROM-CUT1} source-scans
the at most three candidates from one width-three clause, verifies a
singleton strong Krom backdoor, and dry-runs both restrictions.  It accepts
only when every Krom incidence component has encoded size
$N^{1/2+o(1)}$.  The step-capped component enumerator and one-at-a-time
linear-space 2-SAT solver give, on the same machine,
\[
 T(N)\le N^{6/5+o(1)},\qquad S(N)\le N^{1/2+o(1)}.
\]
It decides the earlier sign-balanced lifted-cross family with zero
ordinary-SAT children.  All candidate failures, both assignments,
component discovery, replay and buffers are charged.

This does not promote the ARC to general SAT.  The exact Shannon selector
under $\SAT\in\mathrm P$ uses near-parent SAT answers to choose its route;
inlining preserves a $K$-dependent $N^{K+1+o(1)}$ generic bound, while
exposing the calls gives near-quadratic aggregate invoked volume on the
mixed-window family.  Independent parallel computation segments are
unsound unless their checkpoint states are coupled.  Finally, ordinary
children require source-first certificates that they re-enter the same
covered domain; padding gives accepted PIVOT parents whose unique child is
an uncovered guarded all-fail core.  Thus the favorable point is a promise
point only.  The exact FLMV random-access-machine endpoint remains the
downstream conditional target; P versus NP remains open.
